\documentclass{article}

\usepackage{fancyhdr}
\usepackage{amsmath}
\usepackage{amsthm}
\usepackage{amsfonts}
\usepackage{tikz}
\usepackage[plain]{algorithm}
\usepackage{algpseudocode}
\usepackage{cancel}

\usetikzlibrary{automata,positioning}

\topmargin=-0.45in
\evensidemargin=0in
\oddsidemargin=0in
\textwidth=6.5in
\textheight=9.0in
\headsep=0.25in

\linespread{1.1}

\pagestyle{fancy}
\lhead{\hmwkAuthorName}
\chead{}
\rhead{\hmwkHeaderTitle}
\cfoot{\thepage}

\renewcommand\headrulewidth{0.4pt}
\renewcommand\footrulewidth{0pt}

\setlength\parindent{0pt}

\newcommand{\hmwkTitle}{Lecture\ 6 Practice Exercises}
\newcommand{\hmwkDueDate}{See Canvas for submission time}
\newcommand{\hmwkClass}{Data Structure\ and Abstractions}
\newcommand{\hmwkClassTime}{Section 2}
\newcommand{\hmwkClassInstructor}{Kanat Tangwongsan}
\newcommand{\hmwkAuthorName}{\textbf{Jiraroj Wiruchpongsanon (6781617)}}

\newcommand{\hmwkHeaderTitle}{Data Structure and Abstractions: Lecture 6 Practice}

\title{
    \vspace{2in}
    \textmd{\textbf{\hmwkClass:\ \hmwkTitle}}\\
    \normalsize\vspace{0.1in}\small{Due\ on\ \hmwkDueDate}\\
    \vspace{0.1in}\large{\textit{\hmwkClassInstructor\ \hmwkClassTime}}
    \vspace{3in}
}

\author{\hmwkAuthorName}
\date{}

\renewcommand{\part}[1]{\textbf{\large Part \Alph{partCounter}}\stepcounter{partCounter}\\}

\newcounter{partCounter}
\newcommand{\solution}{\textbf{\large Solution}}

\newtheorem*{theorem}{Theorem}

\theoremstyle{remark}
\newtheorem*{remark}{Remark}

\begin{document}

\maketitle
\newpage

\begin{verbatim}
from copy import copy

def foobar(numbers: List[int]) -> List[int]:
    # make a copy of numbers
    numbers = copy(list(numbers))
    for j in range(1, len(numbers)):
        key = numbers[j]
        i = j-1
        while i >= 0 and numbers[i] > key:
            numbers[i+1] = numbers[i]
            i = i-1
        numbers[i+1] = key
    return numbers
\end{verbatim}

\section*{Problem 1}

Mathematically, what is the running time of the code presented in the puzzle at the beginning of class? We can talk about both worst-case running time and best-case running time. The worst-case running on input of size $n$ is the running time of the hardest input of that size---i.e., one that takes the longest to run. The best-case running time on input of size $n$ is the running time of the easiest input of that size. Your task here is to determine the best-case running time, as well as the worst-case running time, as a function of the input size (the length of the input sequence)?

\medskip
\solution

we let $n = \text{\{length of numbers list\}}$, the `foobar` function is basically an \textbf{Insertion sort} implementation. it (`j`) goes through index $1, 2, 3, ..., (n-1)$, then in each of those iteration, set `key` = {value at index `j` of the list}, and goes through index i = `j` - 1, meanwhile assigning the value of index `i` to index `i` + 1 of the list, and the inner loop only terminates when $i < 0$ or when {the valvue at index `i` of the list} is lower than `key` --- this acts like finding the index that's suitable for our `key`, and lastly, when the inner-loop terminates, set the value of index `i` + 1 to `key` --- which assures that `key` is in the right place, and the list is sorted from the list's head, up to `j`. 

\vspace{1em}

for the \textbf{best case}, the list is already sorted (smallest to greatest), so the inner-loop terminats immediately --- so $\Theta(1)$ for the inner-loop; and we only accounted for the mandatory outer-loop --- of $\Theta(n)$ --- therefore, the total time complexity of the best case is $\boxed{\Theta(n)}$

\vspace{1em}

for the \textbf{worse case}, the list is sorted as well, but in reversed order from what this function intended (greatest to smallest), this will result in the inner-loop running for `j` times, so the total time complexity (already accounted for the outer-loop) would be 

\begin{align*}
    1 + 2 + 3 + ... + n &= \sum_{i=1}^n i \\
    &= \frac{n(n+1)}{2} \in \boxed{\Theta(n^2)}
\end{align*}

\vspace{1em}
\hrulefill

\section*{Problem 2}

Suppose $f(n)$ is $\Theta(s(n))$. Let $g(n) = n \cdot f(n)$. Prove that $g(n)$ is $\Theta(n \cdot s(n))$.

\medskip
\solution

\begin{proof}

\[f(n) \in \Theta(s(n)) \Rightarrow \lim_{n\rightarrow\infty} \frac{f(n)}{s(n)} = c \quad \text{ when } c \ne 0\]

as $g(n) = n \cdot f(n)$ where $f(n) = c \cdot s(n)$

\[
\lim_{n\rightarrow\infty} \frac{g(n)}{n \cdot s(n)}
=
\lim_{n\rightarrow\infty}
\frac{\cancel{n} \cdot c \cdot \cancel{s(n)}}{\cancel{n} \cdot \cancel{s(n)}} = c
\]

therefore

\[g(n) \in \Theta(n \cdot s(n))\]

\end{proof}

\vspace{1em}
\hrulefill

\end{document}
